\clearpage

\section{SMC}
Considérese
\begin{align*}
  \dot x_1&=x_2 \qquad
  \dot x_2=f(x,t)+u,\\
  y&=x_1 \qquad |f(x,t)|\leq L.
\end{align*}
Como
\[
  \dot y=x_2 \qquad
  \ddot y=f(x,t)+u,
\]
el grado relativo es $k=2$. Para
\[
  e=r-y=r-x_1,
\]
se tiene
\begin{align*}
  \dot e=\dot r-x_2 \qquad
  \ddot e=\ddot r-f(x,t)-u.
\end{align*}

Se define
\[
  \boxed{\sigma=\dot e+me} \qquad m>0,
\]
y
\[
  V=\frac{1}{2}\sigma^2.
\]
Entonces
\begin{align*}
  \dot\sigma
  &=m\dot e+\ddot e\\
  &=m\dot r+\ddot r-mx_2-f(x,t)-u.
\end{align*}

Si
\[
  |m\dot r+\ddot r|\leq R,
\]
se propone
\[
  \boxed{u=-mx_2+K\operatorname{sign}(\sigma)}.
\]
Por consiguiente,
\begin{align*}
  \dot V
  &=\sigma\left[m\dot r+\ddot r-f(X,t)
    -K\operatorname{sign}(\sigma)\right]\\
  &\leq -(K-R-L)|\sigma|.
\end{align*}
La condición de convergencia es
\[
  \boxed{K>R+L}.
\]

\subsection{Aplicación al modelo PVTOL}

Para $x$:
\begin{align*}
  e_x&=r_x-x \qquad
  \sigma_x=\dot r_x-\dot x+m_x(r_x-x),\\
  v_x&=-m_x\dot x+K_x\operatorname{sign}(\sigma_x).
\end{align*}

Para $z$:
\begin{align*}
  e_z&=r_z-z \qquad
  \sigma_z=\dot r_z-\dot z+m_z(r_z-z),\\
  v_z&=-m_z\dot z+K_z\operatorname{sign}(\sigma_z).
\end{align*}

Para la actitud:
\begin{align*}
  e_\theta&=r_\theta-\theta \qquad
  \sigma_\theta=\dot r_\theta-\dot{\theta}
  +m_\theta(r_\theta-\theta),\\
  \tau&=-m_\theta\dot{\theta}
  +K_\theta\operatorname{sign}(\sigma_\theta).
\end{align*}

Los parámetros deben satisfacer
\[
  m_x>0 \qquad m_z>0 \qquad m_\theta>0 
\]
\[
  K_x>R_x+L_x \qquad
  K_z>R_z+L_z.
\]

La referencia angular es la señal generada por la linealización:
\[
  r_\theta:=\theta_d=\operatorname{atan2}(-v_x,v_z+g).
\]
La señal $\tau$ es la entrada de control calculada para reducir
$e_\theta$ y conducir $\theta$ hacia $r_\theta$. Como $r_\theta$ puede variar
con el tiempo, si
\[
  \left|
    \ddot r_\theta+m_\theta\dot r_\theta
  \right|\leq R_\theta,
\]
la condición suficiente para el control de actitud es
\[
  K_\theta>R_\theta.
\]